% !TeX program = lualatex
\documentclass[UTF8,zihao=-4,fontset=fandol]{ctexart}
\usepackage[a4paper,left=17mm,right=17mm,top=20mm,bottom=19mm]{geometry}
\usepackage{graphicx,xcolor,tabularray,fancyhdr,dirtree,needspace}
\usepackage[unicode,colorlinks,linkcolor=GuideTeal,urlcolor=GuideTeal]{hyperref}
\IfFileExists{effects.pdf}{\def\GuideEffects{effects.pdf}}{\def\GuideEffects{manual/effects.pdf}}

\definecolor{GuideInk}{HTML}{182E40}
\definecolor{GuideTeal}{HTML}{205B60}
\definecolor{GuideMist}{HTML}{F0F5F6}
\definecolor{GuideBorder}{HTML}{D2DEE1}
\definecolor{GuideGray}{HTML}{425866}

\color{GuideInk}
\setlength{\parindent}{0pt}
\setlength{\parskip}{.35em}
\setlength{\headheight}{15pt}
\pagestyle{fancy}
\fancyhf{}
\fancyhead[L]{\small\sffamily\color{GuideTeal}\textbf{LIQUIDGLASS}}
\fancyhead[R]{\small\sffamily\color{GuideGray}使用说明\enspace /\enspace 2.0.2}
\fancyfoot[R]{\small\color{GuideTeal}\thepage}
\renewcommand\headrulewidth{.3pt}
\ctexset{
  section={format=\Large\bfseries\color{GuideTeal},beforeskip=1.35em,afterskip=.65em},
  subsection={format=\large\bfseries\color{GuideTeal},beforeskip=1.15em,afterskip=.55em}
}
\newcommand\Code[1]{{\ttfamily\small\detokenize{#1}}}
\newcommand\GuideEntry[5]{%
  \begin{minipage}[c]{\linewidth}%
  \setlength{\parskip}{0pt}%
  \textbf{#1}\par
  \Code{#2}\par
  \textcolor{GuideGray}{范围 #3\enspace /\enspace 默认 #4}\par
  #5%
  \end{minipage}%
}
\newcommand\GuideFigure[1]{%
  \raisebox{-.5\height}{\includegraphics[page=#1,width=\linewidth,
    trim=0 2mm 0 1mm,clip]{\GuideEffects}}%
}
\NewTblrEnviron{EffectTable}
\SetTblrInner[EffectTable]{
  colspec={X[9,l,m] X[11,c,m]}, rowhead=1,
  rows={font=\small}, row{1}={bg=GuideTeal,fg=white,font=\bfseries},
  row{even}={bg=GuideMist}, rowsep=3pt, colsep=6pt,
  hline{1,2,Z}={.45pt,GuideTeal}, hlines={.3pt,GuideBorder},
}
\SetTblrOuter[EffectTable]{long,headsep=4pt,footsep=4pt}
\DefTblrTemplate{contfoot-text}{default}{续下页}
\DefTblrTemplate{conthead-text}{default}{续表}
\DefTblrTemplate{caption}{default}{}
\DefTblrTemplate{capcont}{default}{}
\NewTblrEnviron{ReferenceTable}
\SetTblrInner[ReferenceTable]{
  colspec={X[8,l,t] X[12,l,t]}, rowhead=1, rows={font=\small},
  row{1}={bg=GuideTeal,fg=white,font=\bfseries},
  row{even}={bg=GuideMist}, rowsep=4pt, colsep=7pt,
  hlines={.3pt,GuideBorder},
}
\SetTblrOuter[ReferenceTable]{long,headsep=4pt,footsep=4pt,label=none,entry=none}

\begin{document}
\thispagestyle{empty}
\vspace*{5mm}
{\large\sffamily\bfseries\color{GuideTeal} LIQUIDGLASS\hfill 2.0.2}\par
{\fontsize{30pt}{37pt}\selectfont\bfseries 液态玻璃使用说明}\par
{\large\color{GuideGray} 从第一个磨砂盒，到逐项调光和圆角}\par
\vspace{3mm}
{\color{GuideTeal}\rule{\linewidth}{.6pt}}

\section{先做出第一个盒子}
解压后，打开已有的 \Code{example.pdf} 预览。双击 \Code{build.command}，
选择 \Code{example.tex} 重建。成功的 PDF 会出现在对应的 \Code{.tex} 文件旁。
也可以在包目录运行下面的命令。
\begin{verbatim}
./build.command example.tex --no-open
\end{verbatim}

\subsection*{把盒子放进自己的文档}
把以 \Code{liquidglass} 命名的宏包、材质模块和素材目录放在主文档旁，再复制 \Code{lgbuild.py}。
在 macOS 上可一并复制 \Code{build.command}。
下面这段可直接写进普通文档。
\begin{verbatim}
\usepackage{liquidglass}        % 导言区

% 正文
\begin{lgtheorem}{Pythagoras}{pythagoras}
  In a right triangle, $a^2+b^2=c^2$.
\end{lgtheorem}
\end{verbatim}

\subsection*{第一次配置编译环境}
准备 TeX Live 或 MacTeX、Python 3.9 或更新版，以及 Pillow 和 NumPy。
用下列命令在解压目录建立独立的 Python 环境。启动器可自动找到本目录的
\Code{.venv}。采样页面或使用 PDF 底图时还需要 Poppler 的 \Code{pdftoppm}。
\begin{verbatim}
python3 -m venv .venv
.venv/bin/python -m pip install -r liquidglass-requirements.txt
\end{verbatim}
\textcolor{GuideGray}{直接运行一次 LaTeX 可看预览。正式构建会做定位、材质生成和合成，才能看到完整的模糊、柔光与取景效果。}

\clearpage
\section{2.0 光学材质}
TeX 为每个片段记录外轮廓、有效圆角、页坐标与分页状态。构建器由轮廓计算有符号距离场与表面法线，
通过折射近似对底图做逐像素逆向取样。直边、圆角和 Beautybook 的凹圆角因此具有各自的入射方向。

Pillow 负责解码、ICC 转换与 PNG 输出。NumPy 在线性光中计算模糊、薄膜、反光、颗粒和亮部。
\Code{clear} 的边缘保留更多底图细节，\Code{reading} 的正文区保持柔和。分页切口沿盒子的两侧继续，
只有实际首尾才绘制水平边的光学反光。PDF 继续使用 TeX 原生排版与矢量轮廓。

\Code{lg material=reading-lens} 采用银蓝阅读内区、白色方向高光与双层柔影。
圆角的受光位置由表面法线和光源方向决定。折射限于约 2 mm 的受光边缘，底缘延续阅读面。
\Code{lg rim width} 和 \Code{lg rim opacity} 同时调节各边的高光。
\Code{lg rim=hairline} 单独使用矢量细线，可用 \Code{lg rim color} 改色。

用 \Code{lg renderer=legacy} 可在单盒或全局设置中保留旧材质，或给构建器添加 \Code{--renderer legacy}。
\Code{reference} 与 \Code{opaque} 仍采用原版路径。\Code{lg refraction} 在光学预设中按材质换算有效厚度，
相同参数与 1.6.x 可以产生不同的边缘位移。

\clearpage
\section{从类别找到参数}
{\small
\DTsetlength{.25em}{1.25em}{.2em}{.35pt}{.35pt}
\dirtree{%
.1 液态玻璃.
.2 外观.
.3 \hyperref[guide:themes]{主题与预设} · 浅／深、磨砂、清透、原图、柔光.
.3 \hyperref[guide:sources]{背景来源与取景} · 图片、铺图方式、横纵焦点.
.3 \hyperref[guide:material]{材质与柔光} · 模糊、蒙色、颗粒、折射、色彩、对比度与亮部.
.2 轮廓.
.3 \hyperref[guide:shape]{外形与边缘} · 统一半径、单侧、单角、边线与标题签.
.3 \hyperref[guide:lift]{悬浮与接触阴影} · 悬浮预设、宽阴影、贴边阴影.
.2 使用.
.3 \hyperref[guide:interfaces]{接口与构建} · 逐盒、全局、页面、Beautybook 接入、编译选项.
}
}

\medskip
\begin{ReferenceTable}{}
\textbf{常见目标} & \textbf{从这里开始} \\
保留浅银灰纹理 & 背景选 \Code{frost-silver.png}，材质用 \Code{lg reference}。该预设仍保护文字对比度。 \\
让底图更安静 & 少量增加 \Code{lg blur} 或 \Code{lg opacity}。复杂照片可先试 \Code{lg frosted}。 \\
让盒子更轻 & 从 \Code{lg lift=soft} 开始。偏厚重时先减 \Code{lg shadow opacity}。 \\
只改左上角 & 用 \Code{lg radius={3mm,northwest=5pt}}。默认值决定其他三角。 \\
整本使用同样参数 & 在导言区写 \Code{\glasssetup{...}}。单盒的选项可进一步覆盖。 \\
\end{ReferenceTable}

\subsection*{怎样读对照表}
左列给出可复制的选项、范围、默认值和调节方向。右列两盒共用相同的底图、
尺寸与排版，每行只改变该项或一个预设。为了让颗粒、边线或阴影更易辨认，
部分对照提高了固定配套参数。默认值指独立 \Code{glassbox} 使用的基础设置，
Beautybook 的阅读预设见后面的接入说明。
\textcolor{GuideGray}{点击上方类别可直达对应表格。}
\clearpage

\Needspace{82mm}
\section{主题与预设}\label{guide:themes}
先按阅读目的选主题或预设，再微调表中的具体参数。
\begin{EffectTable}[label=none,entry=none]{}
代码、范围与调节方向 & 实编译效果对照 \\
\GuideEntry{lg theme}{lg theme=dark}{light / dark}{light}{一起切换文字颜色和材质色调。两图使用相同底图。} & \GuideFigure{1} \\
\GuideEntry{lg frosted}{lg frosted}{样式关键字}{基础参数近似此预设}{标准磨砂。模糊和蒙色较明显，适合复杂背景。} & \GuideFigure{2} \\
\GuideEntry{lg clear}{lg clear}{样式关键字}{按需启用}{清透玻璃。细节保留较多，最低对比度目标为 4.5:1。} & \GuideFigure{3} \\
\GuideEntry{lg reference}{lg reference}{样式关键字}{按需启用}{适合已处理好的底图。保留全图，仍执行对比度保护。} & \GuideFigure{4} \\
\GuideEntry{lg luminous}{lg luminous}{样式关键字}{按需启用}{加入柔光和悬浮阴影，减弱边缘反光。} & \GuideFigure{5} \\
\end{EffectTable}
\Needspace{82mm}
\section{背景来源与取景}\label{guide:sources}
示例选用可见图形的图片，让裁切方向更容易辨认。来源命令统一列在后面的接口页。
\begin{EffectTable}[label=none,entry=none]{}
代码、范围与调节方向 & 实编译效果对照 \\
\GuideEntry{lg image}{lg image=liquidglass-assets/aurora.png}{路径，或 @page / @default / @folder:目录 / @list:列表}{@page}{只切换当前盒子的底图。相对路径从主文档目录起算。} & \GuideFigure{6} \\
\GuideEntry{lg image fit}{lg image fit=stretch}{cover / stretch}{cover}{等比铺满并裁边，或把整张图拉伸到盒子比例。} & \GuideFigure{7} \\
\GuideEntry{lg focus x}{lg focus x=1}{0–1}{.5}{横向取景从左到右。用窄盒演示横向裁切。} & \GuideFigure{8} \\
\GuideEntry{lg focus y}{lg focus y=1}{0–1}{.5}{纵向取景从上到下。宽盒更容易看到变化。} & \GuideFigure{9} \\
\end{EffectTable}
\Needspace{82mm}
\section{材质与柔光}\label{guide:material}
这些参数处理底图。示例保留同样的正文，便于看背景细节与对比度。
\begin{EffectTable}[label=none,entry=none]{}
代码、范围与调节方向 & 实编译效果对照 \\
\GuideEntry{lg blur}{lg blur=5}{0–30 mm}{2.4}{数值越大，底图上的细节越模糊。正文始终由 TeX 排版。代码只写数字。} & \GuideFigure{10} \\
\GuideEntry{lg opacity}{lg opacity=.78}{0–1}{.30}{提高主题蒙色的比例。浅色更接近浅纸色，深色更接近深蓝灰。} & \GuideFigure{11} \\
\GuideEntry{lg grain}{lg grain=.28}{0–1}{.012}{加入细颗粒。日常可从 .005–.02 试起，图中放大差异。} & \GuideFigure{12} \\
\GuideEntry{lg refraction}{lg refraction=7}{0–10 mm}{.35}{在盒子边缘挤压采样坐标。日常建议从 .2–.7 起试。代码只写数字。} & \GuideFigure{13} \\
\GuideEntry{lg saturation}{lg saturation=1.8}{0–3}{1.08}{0 为灰度，1 保留原饱和度，继续增大则颜色更鲜明。} & \GuideFigure{14} \\
\GuideEntry{lg contrast}{lg contrast=7}{1–21（另受主题上限约束）}{7}{对默认文字色的最低对比度。1 关闭自动修正，实际可达上限见正文。} & \GuideFigure{15} \\
\GuideEntry{lg sheen}{lg sheen=.85}{0–1}{0}{加入局部柔光。深色使用低亮度冷色反光。} & \GuideFigure{16} \\
\GuideEntry{lg sheen size}{lg sheen size=1.1}{0.05–2}{.75}{柔光半径相对于盒子短边的比例。需开启柔光。} & \GuideFigure{17} \\
\GuideEntry{lg sheen x}{lg sheen x=.85}{0–1}{.35}{移动柔光中心。0 在左，1 在右。需开启柔光。} & \GuideFigure{18} \\
\GuideEntry{lg sheen y}{lg sheen y=.94}{0–1}{.22}{移动柔光中心。0 在上，1 在下。需开启柔光。} & \GuideFigure{19} \\
\GuideEntry{lg highlight rolloff}{lg highlight rolloff=1}{0–1}{0}{压低浅色主题的高亮区，减少亮部泛白。深色主题不处理此项。} & \GuideFigure{20} \\
\end{EffectTable}
\Needspace{82mm}
\section{外形与边缘}\label{guide:shape}
方位与单角都写在 lg radius 的大括号内。右图可直接辨认哪一处从圆角切成了直角。
\begin{EffectTable}[label=none,entry=none]{}
代码、范围与调节方向 & 实编译效果对照 \\
\GuideEntry{lg radius}{lg radius=6mm}{0–100 mm，长度须带单位}{glassbox 内弧 4 mm，外弧自动}{统一调整四角。0pt 为直角。} & \GuideFigure{21} \\
\GuideEntry{lg radius 中的 default}{lg radius={default=7mm,northwest=0pt}}{0–100 mm}{取盒子现有设置}{组合里的默认值，也可用 default=。只有未覆盖的角改变。} & \GuideFigure{22} \\
\GuideEntry{lg radius 中的 north}{lg radius={5mm,north=0pt}}{0–100 mm}{沿用统一值}{覆盖上方两角的半径。图右把该处设为直角。} & \GuideFigure{23} \\
\GuideEntry{lg radius 中的 south}{lg radius={5mm,south=0pt}}{0–100 mm}{沿用统一值}{覆盖下方两角的半径。图右把该处设为直角。} & \GuideFigure{24} \\
\GuideEntry{lg radius 中的 west}{lg radius={5mm,west=0pt}}{0–100 mm}{沿用统一值}{覆盖左侧两角的半径。图右把该处设为直角。} & \GuideFigure{25} \\
\GuideEntry{lg radius 中的 east}{lg radius={5mm,east=0pt}}{0–100 mm}{沿用统一值}{覆盖右侧两角的半径。图右把该处设为直角。} & \GuideFigure{26} \\
\GuideEntry{lg radius 中的 northwest}{lg radius={5mm,northwest=0pt}}{0–100 mm}{沿用统一值}{覆盖左上角的半径。图右把该处设为直角。} & \GuideFigure{27} \\
\GuideEntry{lg radius 中的 northeast}{lg radius={5mm,northeast=0pt}}{0–100 mm}{沿用统一值}{覆盖右上角的半径。图右把该处设为直角。} & \GuideFigure{28} \\
\GuideEntry{lg radius 中的 southwest}{lg radius={5mm,southwest=0pt}}{0–100 mm}{沿用统一值}{覆盖左下角的半径。图右把该处设为直角。} & \GuideFigure{29} \\
\GuideEntry{lg radius 中的 southeast}{lg radius={5mm,southeast=0pt}}{0–100 mm}{沿用统一值}{覆盖右下角的半径。图右把该处设为直角。} & \GuideFigure{30} \\
\GuideEntry{lg edge}{lg edge=1}{0–1}{.7}{控制上沿的亮线强度。} & \GuideFigure{31} \\
\GuideEntry{lg edge bottom}{lg edge bottom=1}{0–1}{.24}{控制下沿亮线，以及细内描边的强度。} & \GuideFigure{32} \\
\GuideEntry{lg edge left}{lg edge left=1}{0–1}{.23}{控制左侧亮线的强度。} & \GuideFigure{33} \\
\GuideEntry{lg edge right}{lg edge right=1}{0–1}{.13}{控制右侧暗线的强度。图中用浅青色帮助定位。} & \GuideFigure{34} \\
\GuideEntry{lg edge width}{lg edge width=2.5pt}{0–3 pt}{.65pt}{统一改变边线基准宽度。各边绘制宽度按自身比例缩放。} & \GuideFigure{35} \\
\GuideEntry{lg edge color}{lg edge color=cyan!35!white}{xcolor 颜色}{white}{更换上、下、左亮边及细内描边的颜色。} & \GuideFigure{36} \\
\GuideEntry{lg edge shade color}{lg edge shade color=magenta!35!white}{xcolor 颜色}{LGShadow}{更换右侧暗线用色。需提高右侧强度才容易看出。} & \GuideFigure{37} \\
\GuideEntry{glass tab}{title={Glass}, glass tab}{样式关键字}{按需启用}{需要 title=。将标题移到独立的浅色签条。} & \GuideFigure{38} \\
\end{EffectTable}
\Needspace{82mm}
\section{悬浮与接触阴影}\label{guide:lift}
对照中提高了阴影强度以便观察。实际阅读可以从 soft 预设起试。
\begin{EffectTable}[label=none,entry=none]{}
代码、范围与调节方向 & 实编译效果对照 \\
\GuideEntry{lg lift}{lg lift=cloud}{none / soft / raised / air / cloud}{soft（定理盒为 cloud）}{联动配置宽阴影和接触阴影。图中直接比较原 air 与新版 cloud。} & \GuideFigure{53} \\
\GuideEntry{lg shadow opacity}{lg shadow opacity=.60}{0–1}{.12（soft）}{宽阴影强度。lg shadow 是同一个参数的别名。} & \GuideFigure{40} \\
\GuideEntry{lg shadow blur}{lg shadow blur=7mm}{0–10 mm}{2mm（soft）}{增大后阴影铺得更宽，边缘更软。} & \GuideFigure{41} \\
\GuideEntry{lg shadow xshift}{lg shadow xshift=3mm}{−10 到 10 mm}{0mm（soft）}{宽阴影横移。正值向右。接触层默认跟随，可单独设置。} & \GuideFigure{42} \\
\GuideEntry{lg shadow yshift}{lg shadow yshift=3mm}{−10 到 10 mm}{−.65mm（soft）}{宽阴影纵移。正值向上，负值向下。} & \GuideFigure{43} \\
\GuideEntry{lg shadow spread}{lg shadow spread=3mm}{−2 到 5 mm}{0mm（soft）}{缩小或扩大阴影轮廓的起始范围。} & \GuideFigure{44} \\
\GuideEntry{lg shadow color}{lg shadow color=violet}{xcolor 颜色}{LGShadow}{两层阴影共用的颜色。} & \GuideFigure{45} \\
\GuideEntry{lg shadow steps}{lg shadow steps=48}{4–64 的整数}{24（soft）}{分层越多越平滑，也增加 PDF 绘制量。} & \GuideFigure{46} \\
\GuideEntry{lg contact opacity}{lg contact opacity=.65}{0–1}{0（soft）}{贴边阴影强度，默认关闭。} & \GuideFigure{47} \\
\GuideEntry{lg contact blur}{lg contact blur=4mm}{0–10 mm}{.6mm（soft）}{改变贴边暗线的扩散范围。需开启接触阴影。} & \GuideFigure{48} \\
\GuideEntry{lg contact yshift}{lg contact yshift=2mm}{−10 到 10 mm}{−.3mm（soft）}{接触阴影纵移。正值向上。需开启接触阴影。} & \GuideFigure{49} \\
\GuideEntry{lg contact spread}{lg contact spread=2.5mm}{−2 到 5 mm}{0mm（soft）}{改变接触阴影轮廓的范围。需开启接触阴影。} & \GuideFigure{50} \\
\GuideEntry{lg shadow profile}{lg shadow profile=smooth}{classic / smooth / contour}{classic（float 为 contour）}{图中对照经典与平滑。轮廓渐变在 PDF 透明组内累积衰减，再应用该层的投影强度。} & \GuideFigure{54} \\
\GuideEntry{lg contact xshift}{lg contact xshift=2mm}{−10 到 10 mm}{跟随宽阴影}{只移动接触层。lg contact follows shadow 恢复跟随。} & \GuideFigure{55} \\
\end{EffectTable}
\clearpage
\section{阅读材质与定理预设}
\begin{EffectTable}[label=none,entry=none]{}
代码、范围与调节方向 & 实编译效果对照 \\
\GuideEntry{lg material}{lg material=regular}{regular / clear / opaque / reading / reference}{定理盒为 reading}{复杂图片适合 regular，细节清晰的图片可试 clear。reference 保留底图。} & \GuideFigure{52} \\
\GuideEntry{lg mask}{lg mask=.17}{0–1}{0（reading 为 .17）}{沿整个逻辑盒从左上向右下增加主题蒙色，跨页时保持连贯。} & \GuideFigure{51} \\
\GuideEntry{不透明阅读}{lg material=opaque}{材质预设}{按需启用}{适合浅色输出与复杂背景上的高可读正文。图中和 reading 比较。} & \GuideFigure{56} \\
\GuideEntry{银蓝悬浮阅读}{lg material=reading-lens}{独立可选}{方向高光 .70pt}{正文沿用 reading。细高光随光向渐隐，外侧以双层柔影托起。} & \GuideFigure{57} \\
\end{EffectTable}

\Needspace{150mm}
\subsection{便携卷边定理盒}
\begin{verbatim}
\usepackage{liquidglass}
\glasssettheoremlanguage{zh} % 中文文档先加载 ctex
\begin{lgtheorem}[lg lift=cloud]{勾股定理}{pythagoras}
  直角三角形的三边满足 $a^2+b^2=c^2$。
\end{lgtheorem}
参见定理~\ref{lgthm:pythagoras}。
\end{verbatim}
各环境独立编号，节号变化时重置。下表中的编号环境可加星号。星号形式用方括号传入可选的盒参数，
并用花括号传入标题，标题为空时写 \Code{{}}。证明和解答只接收可选的盒参数。\par
\begin{ReferenceTable}{colspec={X[4,l] X[4,l] X[3,l]},rowsep=2.5pt}
环境 & 标签前缀 & 类别 \\
\Code{lgtheorem} & \Code{lgthm} & 定理 \\
\Code{lgproposition} & \Code{lgprop} & 命题 \\
\Code{lglemma} & \Code{lglem} & 引理 \\
\Code{lgcorollary} & \Code{lgcor} & 推论 \\
\Code{lgdefinition} & \Code{lgdef} & 定义 \\
\Code{lgconjecture} & \Code{lgconj} & 猜想 \\
\Code{lgexample} & \Code{lgex} & 例题 \\
\Code{lgexercise} & \Code{lgexr} & 练习 \\
\Code{lgproblem} & \Code{lgprob} & 问题 \\
\Code{lgremark} & \Code{lgrem} & 注记 \\
\Code{lgapplication} & \Code{lgapp} & 应用 \\
\Code{lgproof} & — & 证明 \\
\Code{lgsolution} & — & 解答 \\
\end{ReferenceTable}
如需换标题签颜色，用 \Code{lg theorem accent=颜色名}。英文为默认语言。原版卷边标题签只出现在首段，
收尾饰带只出现在末段。用 \Code{\begin{lgtimeline}} 和 \Code{\lgtimelineitem} 加入同款时间线。
随包的 \Code{example.tex} 展示“五问建模法”和全部环境，使用 Fandol 字体搭配。已有原工程字体时，
还可在导言区将 \Code{\LiquidGlassFontRoot} 指向原字体目录，路径以斜线结尾。

Apple 的资料将 regular 与 clear 用于不同背景条件，并强调高光、背景扩散和可读性。
本包把这些原则用于静态 PDF，材质取值由 TeX 与构建器实现。
\href{https://developer.apple.com/design/human-interface-guidelines/materials}{Apple Materials}
\quad\href{https://developer.apple.com/videos/play/wwdc2025/219/}{Meet Liquid Glass}
\clearpage
\section{接口、取值联动与构建}\label{guide:interfaces}
\subsection{启用方式与作用范围}
\begin{ReferenceTable}{}
\textbf{接口} & \textbf{用途} \\
\Code{\begin{glassbox}[...]} & 自带可跨页的玻璃盒子。选项里可混用 \Code{title=}、\Code{breakable} 等 tcolorbox 键和表中 \Code{lg ...}。 \\
\Code{liquid glass} & 给已有 tcolorbox 应用玻璃材质。在原样式之后启用，再接逐盒的 \Code{lg ...} 选项。 \\
\Code{\glasssetup{...}} & 为后续玻璃盒追加预设或选项。后写的设置覆盖同名键，逐盒选项优先。 \\
\Code{\glassbackground{...}} & 设后续盒子的通用来源。可传图片路径、\Code{@page} 或 \Code{@default}。 \\
\Code{\glassbackgroundfolder{...}} & 从指定目录按文件名自然顺序逐盒循环。文件名可写成 \Code{01-...}、\Code{02-...}。 \\
\Code{\glassbackgrounds{a.png;b.png}} & 按给定顺序循环，路径之间用分号。建议写项目相对路径。 \\
\Code{\glasspagebackground{...}} & 等比铺满整页，居中裁切。配合 \Code{@page} 时，盒子采样该位置的页底图。 \\
\Code{\clearglasspagebackground} & 清除后续页面的底图。 \\
\Code{LGInk} · \Code{LGPaper} · \Code{LGShadow} & 宏包自带深墨、浅纸和阴影颜色。边线和阴影键也接受其他 xcolor 颜色或色彩混合表达式。 \\
\end{ReferenceTable}


\subsection{选择图片来源}
\begin{ReferenceTable}{}
\textbf{传给 lg image 的值} & \textbf{实际行为} \\
\Code{@page} & 从采集阶段的页面取当前坐标。换页后按各段的实际位置取样。页面图在输出该页时绘制，可在分页之后改下一页的底图。 \\
\Code{@default} & 使用随包的 \Code{liquidglass-assets/default.png}，一张已处理的参考裁片。配 \Code{lg reference} 可直接保留裁片风格。 \\
\Code{@silver} & 使用随包的 \Code{frost-silver.png}，即便携定理与 Beautybook 接入层的银蓝底图。也可传入其他图片的路径。 \\
\Code{@folder:liquidglass-assets/cycle} & 从目录自然排序后循环。与 \Code{\glassbackgroundfolder} 对应。 \\
\Code{@list:a.png;b.png} & 循环清单。和 \Code{\glassbackgrounds} 一样用分号分隔。 \\
\end{ReferenceTable}

可直接使用 PNG、JPEG 等常见位图。PDF 取第一页并需 \Code{pdftoppm}，
其他位图格式取决于 Pillow 的解码支持。指定图片的跨页盒只选图一次，连续铺开后按分页裁切。
圆角与高光只沿首尾的外缘绘制。文件相对主文档目录查找，当前盒的 \Code{lg image} 可覆盖全局来源。
\Needspace{35mm}

\subsection{范围、优先顺序与搭配}
\begin{ReferenceTable}{}
\textbf{搭配} & \textbf{注意点} \\
预设与单项参数 & 写在同一选项列表时，先写 \Code{lg luminous} 等组合样式，再改具体参数。例如 \Code{lg luminous,lg sheen=.25}。 \\
\Code{lg contrast} & 语法允许 1–21。固定默认文字色的可达上限约为浅色 15.39、深色 19.72。一般阅读可从 7 起试。自设文字颜色时需另外目测对比度。 \\
单位 & \Code{lg blur} 和 \Code{lg refraction} 按毫米解释，代码里填裸数字。\Code{lg radius}、\Code{lg edge width}、阴影长度须写 \Code{pt}、\Code{mm} 等单位。 \\
四角覆盖 & 单角覆盖方位，方位覆盖统一值。两个方位同时匹配一个角时，后写的方位优先。圆角组合省略统一值时，未指定处使用现有圆角。极短盒会限制逐角半径。 \\
\Code{lg image fit=cover} & 焦点只有在对应方向确有裁边时才看得出变化。选择 \Code{stretch} 时用整幅图。 \\
\Code{lg sheen} 和 \Code{lg contact opacity} & 对应强度大于 0 时，柔光位置、范围和接触阴影的细调才可见。 \\
\end{ReferenceTable}


\Needspace{78mm}
\subsection{悬浮预设速查}
下表给出预设的关键起始值。开箱默认是 \Code{lg lift=soft}。
\par

\begin{tblr}{
  colspec={Q[l] X[c] X[c] X[c] X[c] X[c]}, rowsep=4pt,
  row{1}={bg=GuideTeal,fg=white,font=\bfseries},row{even}={bg=GuideMist},hlines={.3pt,GuideBorder},
}
预设 & 宽阴影强度 & 模糊 & 纵移 & 层数 & 接触强度 \\
none & 0 & 2 mm & $-.65$ mm & 24 & 0 \\
soft & .12 & 2 mm & $-.65$ mm & 24 & 0 \\
raised & .18 & 3 mm & $-1$ mm & 32 & 0 \\
air & .12 & 5 mm & $-.9$ mm & 40 & .035 \\
cloud & .10 & 6.2 mm & $-1.05$ mm & 40 & .019 \\
\end{tblr}

各预设横移均为 0，阴影色均为 \Code{LGShadow}。air 的宽阴影扩张为 .1\,mm，
cloud 为 $-.35$\,mm，其余为 0。air 的接触模糊和纵移为 .65\,mm、$-.3$\,mm，
cloud 为 1.35\,mm、$-.26$\,mm，其余为 .6\,mm、$-.3$\,mm。各预设接触扩张均为 0。

\Needspace{65mm}
\subsection{Beautybook 接入}
在 Beautybook 主文件导言区，在语言配置之后加载
下面的一行。接入层随正文入口启用。

\smallskip
\Code{\usepackage{stys/beautybook-liquidglass}}\par

它
可作用于书稿里的普通定理盒、代码窗和中式盒。正文外可手动开关。

\begin{ReferenceTable}{}
\textbf{接口} & \textbf{用途} \\
\Code{\BeautyGlassDefault[...]} & 载入 \Code{reading-window} 与 \Code{float} 柔影。\Code{@reading} 在浅色时选银蓝底图，在深色时选照片磨砂底图。 \\
\Code{\BeautyGlassReading[...]} & 载入原银蓝 \Code{reading}、\Code{cloud} 柔影与原有深色主题参数。 \\
\Code{\BeautyGlassSetup{...}} & 在预设之后追加玻璃参数。 \\
\Code{\BeautyGlassBackground{...}} & 在预设之后换源图。 \\
\Code{\BeautyGlassOn} & 打开后续正文盒子的自动接入。\Code{\mainmatter} 会打开。 \\
\Code{\BeautyGlassOff} & 暂停自动接入。前置、后置内容与封面调用会关闭。 \\
\Code{\BeautyGlassStudio[...]} & 旧的阅读预设名称，等同于 \Code{\BeautyGlassDefault}。 \\
\end{ReferenceTable}

中式 \Code{fancybox} 的外框、角饰与底图裁切共用原生半径，可用 \Code{ys/frame/radius} 调整。
独立文档直接使用卷边定理环境，就会载入新版阅读材质和悬浮阴影。

模板的 \Code{macdarkcode} 用白字标题与 \Code{minted} 的原生 Monokai 配色。
代码区减光独立于语法高亮，从标题栏内渐入。

\clearpage
\subsection{构建与项目设置}
\begin{verbatim}
./build.command --list
./build.command example.tex --no-open --sync-next-to-pdf
python3 lgbuild.py example.tex --engine xelatex --dpi 175
\end{verbatim}
\begin{ReferenceTable}{}
\textbf{入口} & \textbf{可调整项} \\
\Code{build.command} & macOS 菜单或指定主文件。\Code{--list} 列文件，\Code{--recursive} 搜子目录，
\Code{--no-open} 只生成 PDF，\Code{--sync-next-to-pdf} 另存同目录的定位文件。
\Code{--keep-window} 或 \Code{--no-pause} 保留终端。 \\
\Code{lgbuild.py} & 可从命令行直接调用，分辨率 \Code{--dpi} 范围 72–300，默认 200。
\Code{--bibliography} 与 \Code{--index} 接受 \Code{auto}、\Code{none}，控制相应辅助工具。 \\
两个入口通用 & \Code{--engine} 取 \Code{auto}、\Code{xelatex}、\Code{lualatex}、\Code{pdflatex}。
\Code{--shell-escape}、\Code{--restricted-shell-escape} 和 \Code{--no-shell-escape} 分别选择完整、受限或关闭外部命令权限。默认关闭。 \\
环境变量 & \Code{LIQUIDGLASS_PYTHON} 可指定启动器使用的 Python。
\Code{LIQUIDGLASS_DRIVER} 可指定启动器使用的构建器路径。
\Code{LIQUIDGLASS_ASSETS} 可为构建器指定内置素材目录。 \\
\end{ReferenceTable}


把 \Code{liquidglass.json} 放在主文档旁即可保存项目设置。根级
\Code{engine} 和 \Code{shell_escape} 是项目默认值，
\Code{documents} 可按主文档文件名覆盖。引擎顺序为显式参数、
文件专用配置、源文件首段的 \Code{!TeX program}、项目默认值，最后采用 LuaLaTeX。
\begin{verbatim}
{
  "engine": "xelatex",
  "shell_escape": false,
  "documents": {
    "other.tex": {"engine": "lualatex"}
  }
}
\end{verbatim}
构建成功后，定位文件也会保存在 \Code{build/}。失败时先看该目录中的同名日志。
如果只看到预览材质，请用启动器重建。图片或工具找不到时，检查文件名、工作环境和路径。
遇到特别高的长盒，降低 \Code{--dpi} 或把内容分为多个逻辑盒。
\end{document}
